\documentclass[aspectratio=169]{beamer}

\usetheme{Madrid}
\usecolortheme{default}

\usepackage{amsmath, amssymb}
\usepackage{graphicx}
\usepackage{booktabs}

\title{Regularized Least Squares for Auto MPG Prediction}
\subtitle{Math 325 Final Project}
\author{Adrian Zhang \and Jacob Pu \and Viktor Liu}
\date{}

\begin{document}

% Slide 1
\begin{frame}
    \titlepage
\end{frame}

% Slide 2
\begin{frame}{Project Goal}

Our project applies ordinary least squares and Tikhonov regularization to the Auto MPG dataset.

\medskip

The main goal is to predict automobile fuel efficiency, measured by miles per gallon (MPG).

\medskip

We compare:

\[
\text{Ordinary Least Squares}
\]

and

\[
\text{Regularized Least Squares / Tikhonov Regularization}
\]

\end{frame}

% Slide 3
\begin{frame}{Dataset}

We use the Auto MPG dataset.

\medskip

The response variable is:

\[
b = \text{MPG}
\]

The predictor variables are:

\[
\text{cylinders, displacement, horsepower, weight, acceleration}
\]

\medskip

In matrix form:

\[
A x \approx b
\]

where \(A\) is the predictor matrix and \(x\) is the coefficient vector.

\end{frame}

% Slide 4
\begin{frame}{Research Questions}

\begin{enumerate}
    \item Which model gives better prediction performance for MPG?
    \item Does regularization improve model stability when predictors are correlated?
    \item How does the choice of \(\lambda\) affect the final result?
\end{enumerate}

\end{frame}

% Slide 5
\begin{frame}{Ordinary Least Squares}

Ordinary least squares minimizes the squared error:

\[
J(x) = \|Ax - b\|^2
\]

The normal equation is:

\[
A^T A \hat{x} = A^T b
\]

If \(A^T A\) is invertible, the solution is:

\[
\hat{x}_{OLS} = (A^T A)^{-1}A^T b
\]

\end{frame}

% Slide 6
\begin{frame}{Tikhonov Regularization}

Tikhonov regularization adds a penalty term:

\[
J(x)=\frac{1}{2}\|Ax-b\|^2+\frac{\lambda}{2}\|x\|^2
\]

The regularized normal equation is:

\[
(A^TA+\lambda I)\hat{x}=A^Tb
\]

So the solution is:

\[
\hat{x}_{ridge}=(A^TA+\lambda I)^{-1}A^Tb
\]

\end{frame}

% Slide 7
\begin{frame}{Why Regularization Helps}

Ordinary least squares only tries to minimize prediction error.

\medskip

Tikhonov regularization balances two goals:

\[
\text{small prediction error}
\]

and

\[
\text{small coefficient size}
\]

\medskip

A larger \(\lambda\) gives stronger regularization.

\medskip

But if \(\lambda\) is too large, the model may underfit the data.

\end{frame}

% Slide 8
\begin{frame}{Implementation}

We used MATLAB to fit both models.

\medskip

Main steps:

\begin{enumerate}
    \item Load the Auto MPG dataset.
    \item Remove missing values.
    \item Use MPG as the response variable.
    \item Standardize the predictor variables.
    \item Split the data into training and testing sets.
    \item Compare OLS and Ridge models using test RMSE.
\end{enumerate}

\end{frame}

% Slide 9
\begin{frame}{RMSE Comparison}

The main performance metric is RMSE:

\[
RMSE = \sqrt{\frac{1}{n}\sum_{i=1}^{n}(b_i-\hat{b}_i)^2}
\]

A lower RMSE means better prediction performance.

\medskip

\begin{table}
\centering
\begin{tabular}{lc}
\toprule
Model & Test RMSE \\
\midrule
OLS & 4.1424 \\
Ridge, \(\lambda=0.001\) & 4.1424 \\
Ridge, \(\lambda=0.01\) & 4.1424 \\
Ridge, \(\lambda=0.1\) & 4.1425 \\
Ridge, \(\lambda=1\) & 4.1428 \\
Ridge, \(\lambda=10\) & 4.1484 \\
Ridge, \(\lambda=100\) & 4.2347 \\
\bottomrule
\end{tabular}
\end{table}

\end{frame}

% Slide 10
\begin{frame}{RMSE Plot from MATLAB}

\begin{center}
    \includegraphics[width=0.78\textwidth]{rmse_plot.png}
\end{center}

\medskip

\centering
{\small Replace \texttt{rmse\_plot.png} with the MATLAB RMSE figure.}

\end{frame}

% Slide 11
\begin{frame}{Coefficient Comparison}

\begin{table}
\centering
\begin{tabular}{lcc}
\toprule
Variable & OLS Coefficient & Best Ridge Coefficient \\
\midrule
Intercept & 23.486 & 23.486 \\
Cylinders & -0.62675 & -0.62674 \\
Displacement & -0.015291 & -0.015426 \\
Horsepower & -2.1755 & -2.1755 \\
Weight & -4.1397 & -4.1396 \\
Acceleration & -0.46348 & -0.46348 \\
\bottomrule
\end{tabular}
\end{table}

\medskip

The best Ridge model uses a very small value of \(\lambda\), so its coefficients are almost the same as OLS.

\end{frame}

% Slide 12
\begin{frame}{Actual MPG vs Predicted MPG}

\begin{center}
    \includegraphics[width=0.72\textwidth]{prediction_plot.png}
\end{center}

\medskip

\centering
{\small Replace \texttt{prediction\_plot.png} with the MATLAB prediction figure.}

\end{frame}

% Slide 13
\begin{frame}{Interpretation of Results}

The OLS model has test RMSE:

\[
RMSE_{OLS}=4.1424
\]

The best Ridge model has:

\[
\lambda = 0.001
\]

\[
RMSE_{ridge}=4.1424
\]

\medskip

So in this train-test split, regularization does not significantly improve prediction accuracy.

\end{frame}

% Slide 14
\begin{frame}{Effect of \(\lambda\)}

When \(\lambda\) is very small:

\[
A^TA+\lambda I \approx A^TA
\]

So Ridge is almost the same as OLS.

\medskip

When \(\lambda\) becomes too large, the penalty term becomes stronger:

\[
\frac{\lambda}{2}\|x\|^2
\]

This shrinks the coefficients too much and can cause underfitting.

\end{frame}

% Slide 15
\begin{frame}{Main Conclusion}

For the Auto MPG dataset:

\begin{itemize}
    \item OLS and Ridge have almost the same best test RMSE.
    \item The best regularization parameter is very small.
    \item Large values of \(\lambda\) increase RMSE.
    \item Tikhonov regularization mainly shows the stability trade-off rather than a large accuracy improvement.
\end{itemize}

\medskip

In summary:

\[
\text{Regularization improves stability, but too much regularization hurts prediction.}
\]

\end{frame}

% Slide 16
\begin{frame}{Questions}

\centering
\Large Thank you.

\medskip

Questions?

\end{frame}

\end{document}
